% Old Chinese Oracle Bones
\section{Old Chinese Oracle Bone Script}

\subsection{Shang Dynasty Divination Records}

Oracle bone inscriptions (c. 1200 BCE) share structural features with Indus script.

\subsubsection{Context}
\begin{itemize}
    \item Administrative and divinatory records
    \item Short inscriptions on durable media
    \item Institutional usage (royal court)
    \item Standardized formulae
\end{itemize}

\subsection{Structural Similarities}

\subsubsection{Inscription Patterns}
\begin{itemize}
    \item Fixed templates (question-procedure-result)
    \item Counting signs (numerals for dates, quantities)
    \item Entity markers (participants, offerings)
    \item Validation markers (royal approval)
\end{itemize}

\subsubsection{Sign Classes}
\begin{itemize}
    \item Numerals: Distinct from semantic signs
    \item Titles/roles: Institutional markers
    \item Offerings: Entity specifications
    \item Results: State outcomes
\end{itemize}

\subsection{Operational Nature}

Oracle bones as operational records:
\begin{itemize}
    \item Document process (divination procedure)
    \item Record state (question, result)
    \item Validate authority (royal participation)
    \item Archive decision (administrative record)
\end{itemize}

\subsection{Differences}

Key distinctions:
\begin{itemize}
    \item Chinese script is clearly phonetic
    \item Oracle bones include narrative elements
    \item More complex syntax
    \item Longer inscriptions average
\end{itemize}

\subsection{Relevance}

Demonstrates that:
\begin{itemize}
    \item Administrative records can follow strict templates
    \item Short inscriptions can be informationally rich
    \item Validation is central to official records
    \item Operational structure compatible with linguistic content
\end{itemize}
